%!TEX root = ../main.tex
% =============================================================================
%  四、问题 1：确定性条件下的计划购电模型      目标 4 页
%  评审在问：确定性版本做对了吗？
%  本节含：① 电量口径与 LP 全式（目标／平衡含弃光松弛／SOC 递推／SOC 界／功率界／日周期）
%          ② 性质引理（充放互斥、能量恒等式）与证明
%          ③ 对偶分析：水价值 ν_t、并网点边际价值 α_t、阈值律与死区宽度
%          ④ 无储能基线 C-01、成本分解恒等式、容量/功率扫描、效率口径极差
%  图名：fig03_q1_dual_price.pdf / fig10_strategy_mechanism.pdf
%  单位纪律：功率 kW、电量 kWh、电价 元/kWh、费用 元；Δt=1/6 h；P̄Δt=833.3333 kWh
% =============================================================================

\section{问题 1：典型日确定性计划}\label{chap:q1}

\subsection{模型建立}

\subsubsection{模型与电量口径}

\textbf{口径。} 采用电量口径：决策变量均以电量（kWh）为单位，电价 \(p_t\)（元/kWh）与电量相乘即
费用，\textbf{不再乘 \(\Delta t\)}；负载与光伏（kW）进模型前乘以 \(\Delta t\) 化为电量。
电价与负载逐日相同、光伏给出一日预测，故只需求解一个典型日（附件 1），信息结构为“当日真值已知”。
储能参数：额定容量 \(12\,000\)~kWh，运行区间 \([\underline S,\overline S]=[1200,10\,800]\)~kWh
（可用摆幅 9600~kWh），额定功率 \(\overline P=5000\)~kW，单边效率 \(\eta_c=\eta_d=0.9\)（假设 A2）。

决策变量为 \((E_t,c_t,d_t,s_t,S_t)\)，\(t\in\mathcal T\)，\(|\mathcal T|=144\)。模型为

\begin{subequations}\label{eq:q1-model}
\begin{align}
  \min_{E,c,d,s,S}\quad & C^{\text{plan}}=\sum_{t}p_tE_t \label{eq:q1-obj}\\
  \text{s.t.}\quad
  & E_t+(PV_t-L_t)\Delta t+d_t-c_t-s_t=0, && \forall t \label{eq:q1-balance}\\
  & S_t=S_{t-1}+\eta_c c_t-d_t/\eta_d, && \forall t \label{eq:q1-soc}\\
  & \underline S\le S_t\le\overline S,\quad 0\le c_t,d_t\le \overline P\Delta t, && \forall t \label{eq:q1-socbound}\\
  & E_t\ge0,\quad s_t\ge0, && \forall t \label{eq:q1-nonneg}\\
  & S_0=S_{144}=S^{init}=6000\ \text{kWh}. \label{eq:q1-periodic}
\end{align}
\end{subequations}

式~\eqref{eq:q1-balance} 即“微网供给不低于小区负载”的代数表达，\((PV_t-L_t)\Delta t\) 为净负荷电量
（盈余为负）。光伏盈余可能超过充电功率与负载之和，故以\textbf{等式加松弛 \(s_t\ge0\)} 写出以保证
可行；典型日的日内盈余低于摆幅 9600~kWh 与单槽吸纳能力 \(\overline P\Delta t=833.3333\)~kWh，
故最优 \(s_t\equiv0\)。同一典型日上“含松弛等式／不含松弛等式／不等式”三种写法的
最优值完全相同（均为 \(35\,126.9486\)~元），保留 \(s_t\) 的意义在于实际日期序列
（窗口内光伏峰值 \(10\,216.2000\)~kW）的可行性。式~\eqref{eq:q1-periodic} 即题面“两端储电量相同”，
本文进一步取两者均等于附件给定的初值 6000~kWh。

模型规模为 \(4\times144=576\) 个变量、288 条约束（显式保留 \(s_t\) 时为 720 个变量），是连续 LP；
充放互斥由模型结构自动成立，\textbf{无需 0--1 变量}。

\subsubsection{两条结构性质}

\begin{proposition}[充放互斥]\label{prop:no-simul}
若 \(\eta_c\eta_d<1\)，则最优解中不存在 \(c_t,d_t\) 同时为正的时段。
\end{proposition}

\begin{proof}
设某 \(t\) 有 \(c_t>0,d_t>0\)。取 \(\delta=\min\{c_t,\ d_t/(\eta_c\eta_d)\}>0\)，作
\(c_t'=c_t-\delta\)、\(d_t'=d_t-\eta_c\eta_d\delta\ge0\)。由式~\eqref{eq:q1-soc}，
储电量变化 \(\eta_c(-\delta)+\eta_c\eta_d\delta/\eta_d=0\)，故储电量轨迹与全部界约束不变；
由式~\eqref{eq:q1-balance}，为保持平衡购电量必须变为
\(E_t'=E_t-\delta(1-\eta_c\eta_d)\)。\textbf{这一步需要分两种情形}，因为
\(E_t'<E_t\) 可能触及 \(E_t\ge0\) 的下界（式~\eqref{eq:q1-nonneg}）：

\emph{情形一（\(E_t\) 有余量）}：若 \(E_t\ge\delta(1-\eta_c\eta_d)\)，则 \(E_t'\ge0\)，
盈余处置量 \(s_t\) 不变，新解可行且目标下降
\(p_t\delta(1-\eta_c\eta_d)>0\)，与原解最优矛盾。

\emph{情形二（\(E_t\) 贴界）}：若 \(E_t<\delta(1-\eta_c\eta_d)\)，则保持 \(E_t'=E_t\)
并把松弛量增加 \(\delta(1-\eta_c\eta_d)\)，即 \(s_t'=s_t+\delta(1-\eta_c\eta_d)\ge0\)
（\(s_t\) 只有下界 \(0\)、无上界），储电量轨迹与全部界约束同样不变。此时目标\emph{不变}，
新解与原解同为最优——即\textbf{不能排除“存在同时充放的最优解”}，只能断言
“存在一个同时充放为零的最优解”（对 \(c_t,d_t\) 作上述变换即可构造）。
因此命题的严格表述应加一句：\emph{最优值可由不含同时充放的解达到}；本文结果文件给出的
是这一支解，式~\eqref{eq:q1-nonneg} 的 \(E_t\ge0\) 与目标系数 \(p_t>0\) 保证了
“贴界槽”上同时充放不产生任何收益。
\end{proof}

命题~\ref{prop:no-simul} 只用 \(\eta_c\eta_d<1\)、\(p_t>0\) 与平衡式结构，对问题 1--4 的全部变体
一致成立；\textbf{边界情形}（\(E_t=0\) 且 \(s_t=0\)，典型地出现在光伏盈余槽）由上述情形二覆盖：
此时变换后的解与原解同值，故互斥性只能按“最优值可由互斥解达到”理解，而不能按“每个最优解都互斥”理解。
数值验证：结果文件中同时充放时段数 0，储电量递推残差 \(9.09\times10^{-13}\)~kWh，
充放比偏差 \(4.44\times10^{-16}\)（审计项 AU5）。

\begin{lemma}[能量恒等式]\label{lem:energy}
逐槽由式~\eqref{eq:q1-balance} 解出 \(E_t\) 并对 \(t\) 求和，再用式~\eqref{eq:q1-soc} 消去充放电量，得
\begin{equation}
  \sum_{t}E_t=\sum_{t}(L_t-PV_t)\Delta t+\sum_{t}c_t-\sum_{t}d_t
  \ \xrightarrow{\ S_{144}=S_0\ }\
  \sum_{t}(L_t-PV_t)\Delta t+\bigl(1-\eta_c\eta_d\bigr)\sum_{t}c_t .
  \label{eq:q1-energy}
\end{equation}
\end{lemma}

引理~\ref{lem:energy} 有三点用途。① \textbf{它是“恒价无套利”的代数根源}：若日内电价恒为
\(\bar p\)，则目标 \(\bar p\sum_tE_t\) 只随循环量 \(\sum_tc_t\) 变化，与充放时序无关，储能\emph{没有}
时段套利价值；反之，题面背景段所指的“低充高放”只有在日内电价随时刻变化时才有对象，这从代数上支持
假设 A11（附件 1 的 144 点曲线逐日重复）。② \textbf{它是零成本的可行性回归校验}：逐槽核验
式~\eqref{eq:q1-energy}，典型日两侧之差在 \(10^{-4}\)~kWh 量级（即四位小数取整量级）。
③ \textbf{它分离了两问的价值来源}：问题 1 的节省来自改变 \(\sum_tE_t\) 的\emph{时段结构}（套利），
问题 2 的安全边际来自改变其\emph{分布}（降低高价缺口的期望），两者可叠加但不可互替。

\subsubsection{对偶分析与阈值型最优策略}

记 \(\nu_t\) 为式~\eqref{eq:q1-soc} 的影子价格，即\textbf{单位储电量的边际价值}（水价值，元/kWh）；
\(\alpha_t\) 为式~\eqref{eq:q1-balance} 的影子价格（并网点电量的边际价值，元/kWh，取非负约定）。
在储电量界与功率界均不活跃时，对 \(c_t,d_t\) 的平稳性条件给出互补松弛对

\begin{equation}
  z^{c}_t=\alpha_t-\eta_c\nu_t\ge0,\quad c_tz^{c}_t=0;\qquad
  z^{d}_t=-\alpha_t+\nu_t/\eta_d\ge0,\quad d_tz^{d}_t=0 .
  \label{eq:q1-kkt}
\end{equation}

其经济含义为：多充 1~kWh 的边际代价是 \(\alpha_t\)，换回 \(\eta_c\)~kWh 储电、价值
\(\eta_c\nu_t\)，故仅当 \(\alpha_t\le\eta_c\nu_t\) 时充电；多放 1~kWh 可替代 \(\alpha_t\) 的购电成本，
代价是储电量减少 \(1/\eta_d\)~kWh，故仅当 \(\alpha_t\ge\nu_t/\eta_d\) 时放电。当 \(E_t>0\) 时由互补
松弛得 \(\alpha_t=p_t\)（实测偏差 \(\le10^{-6}\)）；当 \(E_t=0\) 时 \(\alpha_t\le p_t\)。由此得

\begin{proposition}[阈值律]\label{prop:threshold}
在储电量界不活跃（\(\underline S<S_t<\overline S\)）的时段，最优充放策略由两条阈值刻画：
\(\alpha_t\le\eta_c\nu_t\) 时充电，\(\alpha_t\ge\nu_t/\eta_d\) 时放电，而
\begin{equation}
  \eta_c\nu_t\ <\ \alpha_t\ <\ \nu_t/\eta_d
  \label{eq:q1-deadband}
\end{equation}
时 \(c_t=d_t=0\)，构成“死区”。
\end{proposition}

死区宽度随往返效率下降而变宽：两条阈值之比为 \((\nu_t/\eta_d)/(\eta_c\nu_t)=1/(\eta_c\eta_d)=1.2346\)。
即\textbf{放电时段的并网点电量边际价值须高于充电时段的 1.2346 倍，跨时段套利才有利可图}；在购电
时段该判据化为放电电价高于充电电价 1.2346 倍，此即式~\eqref{eq:ch1-arb} 的盈亏平衡比。

\textbf{适用边界。} 命题~\ref{prop:threshold} 在“纯套利子问题（变量仅 \(c_t,d_t,S_t\)）且储电量界
不活跃“时精确成立；界活跃时须回到式~\eqref{eq:q1-kkt} 并计入界乘子；功率界活跃时 \(z^c_t\) 或
\(z^d_t\) 不再为零，简单阈值关系失效。数值验证分两层：144 个时段上按 \(\alpha_t\) 判定充／放／死区的
\textbf{违反次数为 0}；界不活跃的 97 个时段子集（充电槽 42、放电槽 36）上违反次数同样为 0，
两条互补松弛恒等式最大残差为 0。实测 \(\nu_t\in[0.4711,1.1305]\)、\(\alpha_t\in[0.3713,1.2561]\)~元/kWh
均\textbf{时变}，故不存在“全天统一的电价阈值”；正确表述是两级结构——先由 LP 对偶确定时变水价值曲线，
再逐时段比较 \(\alpha_t\) 与两条阈值 \evid{C-Q1-02}。

\textbf{三处计数必须能对上账}（否则读者无法判断“违反 0 次”的分母是什么）：图~\ref{fig:dualprice} 标注
全天\textbf{充电槽 45、放电槽 37、空闲槽 62}（合计 \(144\)）；其中储电量界\textbf{不活跃}的子集为
\textbf{97} 槽 \(=42\)（充电）\(+36\)（放电）\(+19\)（空闲），即 45 个充电槽中有 3 个、37 个放电槽中
有 1 个被储电量界或功率界顶住；附录 \appref{tab:scan-app}{D.3 容量与功率扫描明细}另报界活跃计数
\(19/0/21/8\)（充电功率界／放电功率界／储电量上界／下界），是\textbf{按约束出现次数}统计（同一槽可同时
顶住两条界），故与“不活跃 97 槽”不构成简单相加关系。阈值律的一致性结论只在\textbf{第二组（97 槽）}
上精确成立，第一组（45 槽）中含界活跃槽，须计入界乘子。

\begin{figure}[htbp]
  \centering
  \includefig[width=0.60\linewidth]{fig03_q1_dual_price.pdf}
  \caption{问题 1 的阈值律与充放电量（上子图：电价 \(p_t\)、储电边际价值 \(\nu_t\) 与两条阈值
            \(\eta_c\nu_t\)、\(\nu_t/\eta_d\) 构成的死区与充电窗口；下子图：最优充/放电量，
            充电为正、放电为负；对象：附件 1 典型日的 144 个 10~min 槽；单位：元/kWh、kWh）}
  \label{fig:dualprice}
\end{figure}

\subsubsection{无储能基线与成本分解}

无储能时 \(c_t=d_t=s_t=0\)，由式~\eqref{eq:q1-balance} 得闭式
\(E_t=\Delta t\max(L_t-PV_t,0)\)，据此定义基线

\begin{equation}
  C^{base}=\sum_t p_t\Delta t\max\bigl(L_t-PV_t,0\bigr),\qquad
  \sum_t p_tE_t=C^{base}+C^{batt},
  \label{eq:q1-decomp}
\end{equation}

其中 \(C^{batt}=\sum_t p_t\bigl(\Delta t\min(L_t-PV_t,0)+c_t-d_t+s_t\bigr)\) 为储能与弃光带来的成本
增量（典型日取负值即费用下降）。式~\eqref{eq:q1-decomp} 是\textbf{逐项恒等式}而非两个独立最优值
之和，\(C^{batt}\) 的最小化仍需求解式~\eqref{eq:q1-model}。基线有三个变体，\textbf{不得混用}
（表~\ref{tab:cbases}），否则全文多个“节省率”互不可比。

\begin{table}[htbp]
  \centering
  \caption{无储能基线 \(C^{base}\) 的三个变体及其适用问题（单位：元）}
  \label{tab:cbases}
  \footnotesize
  \begin{tabularx}{\linewidth}{C{1.0cm} L{3.6cm} L{3.0cm} C{1.6cm} Y}
    \toprule
    变体 & 口径 & 电价来源 & 适用问题 & 数值 \\
    \midrule
    V1 & 附件 1 典型日，单日周期 & 附件 1（固定价） & 问题 1 & 48\,052.0466 \\
    V2 & 334 天窗口，逐日独立 & 附件 1（固定价） & 问题 2、3 & 16\,407\,319.63 \\
    V3 & 334 天窗口，逐日独立 & 附件 4（波动价） & 问题 4 & 17\,026\,782.86 \\
    \bottomrule
  \end{tabularx}
\end{table}

\subsection{模型求解与结果分析}

\subsubsection{求解结果与题面格式表}

在附件 1 典型日上求解式~\eqref{eq:q1-model}：全天购电量 \(59\,482.6990\)~kWh（\calref{C-02}）、
购电费 \(35\,126.9486\)~元，相对 \(C^{base}=48\,052.0466\)~元（变体 V1）节省
\(12\,925.0980\)~元，\textbf{节省率 \(26.90\%\)}（\calref{C-04}）\evid{C-Q1-03}；
充电量 \(20\,740.6661\)~kWh、放电量 \(16\,799.9396\)~kWh，二式互为倒数——
\(c/d=1.2346=1/(\eta_c\eta_d)\)，其倒数 \(d/c=0.8100=\eta_c\eta_d\) 即往返效率本身（\calref{C-07}）；
最优弃光 \(0.0000\)~kWh（\calref{C-06}）\evid{C-Q1-01}；储电量全程位于
\([1200.0000,10\,800.0000]\)~kWh，两端均为 \(6000.0000\)~kWh（严格日周期）。由
式~\eqref{eq:q1-decomp} 代入上述数值得 \(C^{batt}=-12\,925.0980\)~元，与目标函数直接相减一致。
表~\ref{tab:q1-buy}、表~\ref{tab:q1-soc} 按题面表 1、表 2 格式给出指定时段与 4~h 块的结果
（与 \texttt{result1.xlsx} 逐格一致）。

\begin{table}[htbp]
  \centering
  \caption{微网在指定时间段的购电量及全天的购电量和购电费（问题 1，单位：\unit{kWh} / \unit{元}）}
  \label{tab:q1-buy}
  \small
  \begin{tabular}{cccccc}
    \toprule
    时间段 & 购电量 & 时间段 & 购电量 & 时间段 & 购电量 \\
    \midrule
    \(10{:}00\)--\(10{:}10\) & 0.0000 & \(12{:}00\)--\(12{:}10\) & 480.4124 & \(14{:}00\)--\(14{:}10\) & 0.0000 \\
    \(16{:}00\)--\(16{:}10\) & 445.4317 & \(18{:}00\)--\(18{:}10\) & 531.8940 & \(20{:}00\)--\(20{:}10\) & 0.0000 \\
    \midrule
    \multicolumn{2}{c}{全天购电量（\unit{kWh}）：59\,482.6990} &
    \multicolumn{4}{c}{全天购电费（\unit{元}）：35\,126.9486} \\
    \bottomrule
  \end{tabular}
  \\[1ex]
  {\footnotesize 表注：时段按题面表述取区间左端时刻；\texttt{result1.xlsx} 沿用模板既有的时段末端
  标签，其数值与本文槽位按\textbf{位置}一一对应（槽 \(j\) 写入第 \(j+1\) 个数据行），标签仅作显示。}
\end{table}

\begin{table}[htbp]
  \centering
  \caption{储能设备在指定时间段的充放电量及 \(0{:}00\) 和 \(24{:}00\) 的储电量（问题 1，单位：\unit{kWh}）}
  \label{tab:q1-soc}
  \small
  \begin{tabular}{cccccc}
    \toprule
    时间段 & 充电量 & 放电量 & 时间段 & 充电量 & 放电量 \\
    \midrule
    \(0{:}00\)--\(4{:}00\)   & 4500.0000 & 0.0000    & \(4{:}00\)--\(8{:}00\)   & 833.3333  & 6365.8412 \\
    \(8{:}00\)--\(12{:}00\)  & 4787.9643 & 1702.9970 & \(12{:}00\)--\(16{:}00\) & 5286.0352 & 91.1014 \\
    \(16{:}00\)--\(20{:}00\) & 0.0000    & 5780.1319 & \(20{:}00\)--\(24{:}00\) & 5333.3333 & 2859.8681 \\
    \midrule
    \multicolumn{2}{c}{\(0{:}00\) 储电量：6000.0000} &
    \multicolumn{4}{c}{\(24{:}00\) 储电量：6000.0000} \\
    \bottomrule
  \end{tabular}
  \\[1ex]
  {\footnotesize 表注：每段含 24 个 10~min 槽；6 段之和与逐槽求和一致——充电 20\,740.6661~kWh、
  放电 16\,799.9396~kWh（\calref{C-07}）。}
\end{table}

\begin{figure}[htbp]
  \centering
  \includefig[width=0.54\linewidth]{fig10_strategy_mechanism.pdf}
  \caption{最优策略的机制解释（三子图共享时间轴：光伏与负载功率（\unit{kW}）、
            电价与储电量（\unit{元/kWh}、\unit{kWh}，储电量右轴标出上下界 1200/10800~\unit{kWh}）、
            充放电量（\unit{kWh}）；对象为附件 1 典型日的 144 个 10~\unit{min} 时段，
            竖直虚线连接光伏盈余时段、充电时刻与低价时段）}
  \label{fig:mechanism}
\end{figure}

图~\ref{fig:mechanism} 显示充电时段与光伏盈余时段、低价时段对齐，放电时段与高价时段对齐，
两组时段只\textbf{部分重合}，说明主机制是面向价格的跨时段套利，而非“把午间盈余搬到傍晚”。

\subsubsection{策略机制与对偶解释}

最优解的时段分布给出机制的量化画像：充电量 \(20\,740.67\)~kWh 落于低谷
\(0{:}00\)--\(4{:}00\)（\(4\,500.00\)~kWh）、\(20{:}00\)--\(24{:}00\)（\(5\,333.33\)~kWh）
与午间 \(12{:}00\)--\(16{:}00\)（\(5\,286.04\)~kWh）；放电量 \(16\,799.94\)~kWh 集中于
\(4{:}00\)--\(8{:}00\)（\(6\,365.84\)~kWh）与 \(16{:}00\)--\(20{:}00\)（\(5\,780.13\)~kWh）。
两段低谷占全日充电量的 \(47.4\%\)、午间盈余段仅占 \(25.5\%\)，放电在两段高峰合计占 \(72.3\%\)——
故主机制是\textbf{面向电价的跨时段搬运}，而非“把午间盈余存到傍晚”。

\textbf{对偶解释。} 对偶量给出可解释的策略形式（图~\ref{fig:dualprice}）：实测
\(\nu_t\in[0.4711,1.1305]\)、\(\alpha_t\in[0.3713,1.2561]\)~元/kWh，\(E_t>0\) 时
\(\alpha_t=p_t\)（偏差 \(\le10^{-6}\)），144 个槽的阈值判定\textbf{违反次数为 0}，死区比
\(1/(\eta_c\eta_d)=1.2346\) 与实测充放比一致；储电量日内触及下限 \(1200.00\) 与上限
\(10\,800.00\)~kWh，两端均为 \(6000.00\)~kWh。\evid{C-Q1-02}

\subsection{检验与灵敏度分析}

\subsubsection{内部一致性与基线对比}

\textbf{内部一致性。} 三条恒等式全部成立：供需平衡残差 \(1.49\times10^{-13}\)~kWh、储电量递推残差
\(9.09\times10^{-13}\)~kWh、日周期偏差 \(\le10^{-12}\)~kWh；储电量越界、同时充放与负购电的时段数
\textbf{均为 0}；费用相对独立复算误差 \(\le10^{-6}\)（阈值 \(10^{-4}\)）。阈值律一致性的实测为：
144 个槽的阈值判定违反次数 0、互补松弛最大残差 \(\le10^{-12}\)。\textbf{该检验的作用边界须写明}：
\(\nu_t,\alpha_t\) 本身就是求解器返回的对偶量，互补松弛是任何 LP 最优解对其对偶解的自动性质，
故它只说明求解器返回的解与其自带对偶量自洽、且本文的对偶映射没写错，\emph{不构成}模型正确性的
独立证明（完整讨论见 \appref{sec:verify-three-app}{D.1 独立审计报告}）。

\textbf{基线对比。} 不做优化（无储能）时全天购电费 \(C^{base}=48\,052.0466\) 元（变体 V1），
优化后为 \(35\,126.9486\) 元，\textbf{下降 \(12\,925.0980\) 元（\(26.90\%\)）}（\calref{C-04}）。
按引理~\ref{lem:energy}，同价下 \(\sum_tE_t\) 只随循环量变化，故这笔下降与电价曲线的时间结构一一对应，
来自储能对电量的跨时段搬运与弃光消纳，而非“容量更大所以买得更少”。

\subsubsection{容量、功率与效率口径的灵敏度分析}\label{sec:q1-scan}

在保持 10\%--90\% 放电深度（\(\underline S=0.1E_{cap}\)、\(\overline S=0.9E_{cap}\) 同步移动）下扫描
额定容量，并单独扫描功率上限（完整档位见 \appref{tab:scan-app}{D.3 容量与功率扫描明细}）：

\begin{equation}
  \begin{aligned}
    E_{cap}:12\,000\to28\,800\ \text{kWh} &\Rightarrow \Delta C=-6272.42\ \text{元}\ (-17.86\%),\\
    \overline P:5000\to10\,000\ \text{kW} &\Rightarrow \Delta C=-166.36\ \text{元}\ (-0.47\%),\\
    \text{仅解除 }\overline S &\Rightarrow -22.14\%\ (27\,351.55\ \text{元}),\qquad
    \text{仅解除 }\overline P &\Rightarrow -1.66\%\ (34\,544.43\ \text{元}).
  \end{aligned}
  \label{eq:q1-scan}
\end{equation}

界活跃统计：充电功率界在 \(19/144\) 个时段被打满，放电功率界 \(0/144\)；储电量上界活跃
\(21/144\)、下界活跃 \(8/144\)。故\textbf{容量摆幅是主要瓶颈，充放电功率仅在部分时段同时构成约束}，
两者量级相差一个数量级；这与 \appref{sec:sens-app}{D.3 单参数灵敏度完整档位} 中储能参数双向扫描的排序方向一致
（Sobol 全局灵敏度另显示，全天购电费量级主要由电价水平决定；详见
\appref{tab:sobol-app}{D.3 Sobol 全局灵敏度指数}）。

\textbf{效率口径。} 题面“充放电效率为 90\%”未区分充、放方向，故须显式声明口径
（主口径的选择理由见假设 A2）。以 \(\eta_c,\eta_d\) 的四种读法重解，得节省额

\begin{equation}
  \Delta C:\quad
  \underbrace{12\,925.10}_{0.9/0.9\ (\text{主口径})},\quad
  \underbrace{14\,250.99}_{\text{往返 }0.9},\quad
  \underbrace{14\,635.52}_{\text{仅充 }0.9},\quad
  \underbrace{13\,883.44}_{\text{仅放 }0.9}\quad(\text{元}),
  \label{eq:q1-eff}
\end{equation}

极差 \(1710.42\) 元，占主口径节省额的 \(13.2\%\)。故效率口径必须声明，否则同一模型可报出相差
八分之一的节省率。

把本问的单参数扫描按“结论是否改变”归纳为表~\ref{tab:sens-q1}，完整档位实测值见 \appref{sec:sens-app}{D.3 单参数灵敏度完整档位}；
全局灵敏度（Sobol 指数）回答的是另一个问题——“哪个参数支配输出方差”，见下段。

\begin{table}[htbp]
  \centering
  \caption{问题 1 的参数灵敏度：扰动档位与“结论是否改变”的判断（费用单位 元，电量单位 kWh；
           单参数扫描除终端条件外均在同一组典型日参数下求解；完整档位实测值见
           \appref{sec:sens-app}{D.3 单参数灵敏度完整档位}）}
  \label{tab:sens-q1}
  \small
  \setlength{\tabcolsep}{4pt}
  \begin{tabularx}{\linewidth}{L{2.05cm} L{3.15cm} Y}
    \toprule
    参数 & 扰动档位 & 结论是否改变（含反转点） \\
    \midrule
    效率口径 \(\eta_c,\eta_d\) & 单边各 0.9（\textbf{主口径}）、往返 0.9 等效单边、
        仅充电侧 0.9、仅放电侧 0.9
      & 区间内结论不变（四种读法下储能均有利可图），但节省额极差 \textbf{13.2\%}
        （式~\eqref{eq:q1-eff}）——同一模型可因效率口径不同报出相差八分之一的数 \\
    储电量运行区间 & 保持 10\%--90\% DOD：容量 7\,200--48\,000~kWh 共 7 档
      & 区间内结论不变：容量摆幅始终是主要瓶颈，仅解除容量上限可再省 22.14\% \\
    充放功率上限 \(\overline P\) & 1\,250、2\,500、5\,000（基准）、7\,500、10\,000~kW
      & 区间内结论不变：功率在 5000~kW 以上边际很小（\(\le0.47\%\)），
        但收紧到 2\,500~kW 时节省率由 26.90\% 降到 24.27\%、1\,250~kW 时降到 16.41\% \\
    终端条件 \calref{C-08} & 窗口：\(S_T\ge S_0\)（主口径）／自由终端；典型日：日周期／自由环状
      & 三组差异均不超过 0.12\%，结论不变；窗口口径下自由终端反而略贵（\(+0.08\%\)），
        说明终端约束并非靠“掏空”储能省钱 \\
    \bottomrule
  \end{tabularx}
\end{table}

Sobol 全局灵敏度须\textbf{区分输出量}，并\textbf{同时读一阶指数 \(S_1\) 与总效应指数 \(S_T\)}
（完整表见 \appref{tab:sobol-app}{D.3 Sobol 全局灵敏度指数}）。对\textbf{典型日全天购电费}：电价水平
\(S_1=0.7588\)、\(S_T=0.7758\) 主导，其次是负载水平（\(0.0765/0.0911\)）、储电量上限（\(0.0130/0.0158\)）
与光伏水平（\(0.0103/0.0291\)）；8 个参数的一阶指数\textbf{合计 \(0.8765\)}、总效应\textbf{合计
\(0.9345\)}，即\textbf{交互项只贡献约 5.8 个百分点}（主要来自电价水平与负载水平的交互），
不存在“强交互主导”的情形。对\textbf{节省率}：负载水平（\(0.2157/0.2257\)）与电价水平（\(0.1821/0.2000\)）
主导，一阶合计约 \(0.5093\)、总效应合计 \(0.5527\)，交互约 4.3 个百分点；储能与效率参数的一阶指数
均不超过 \(0.007\)。\textbf{读数边界（重要）}：购电费近似“价格 \(\times\) 电量”的乘性结构使电价水平的
一阶指数\textbf{接近平凡}——它按比例缩放整个费用，故必然主导方差；因此该结论只回答“哪个参数支配
\emph{输出方差}“，\textbf{不能}读成”电价水平决定购电费\emph{量级}“这一均值效应命题（方差分解看不见
均值），也不能反推“储能参数不重要”（它们改变的是可套利电量，而不是费用的尺度）。

\subsubsection{稳健性小结}

本问结论在扰动下稳定，但有一处量级敏感性须声明。① \textbf{方向不变}：效率四种读法、终端条件三组
对照、容量 7 档与功率 5 档扫描下，储能均降低购电费。② \textbf{瓶颈排序不变}：容量摆幅始终是主要
瓶颈（仅解除容量上限可再省 \(22.14\%\)），充放功率在 \(5000\)~kW 以上边际 \(\le0.47\%\)；
但把功率上限收紧到 \(2500\)~kW 时节省率由 \(26.90\%\) 降到 \(24.27\%\)、收紧到 \(1250\)~kW 时降到
\(16.41\%\)，即\textbf{功率上限低于题面值才会改变结论}。③ \textbf{一处口径敏感性}：四种效率读法下
节省额在 \(12\,925.10\)--\(14\,635.52\) 元之间，极差 \(1710.42\) 元（占主口径的 \(13.2\%\)），
故主口径（单边各 \(0.9\)）必须在假设 A2 中声明并在全部结果表中固定使用；
终端条件的影响则不超过 \(0.12\%\)，不构成结论风险。
